\documentclass[10pt]{extarticle}
\usepackage{apMod}
\usepackage{hyperref}
% Licensed: https://www.gnu.org/licenses/gpl-3.0.html
% Template author: Anindya Prithvi
\vspace{-20pt}
% === Personal Information ===
% === Personal Information ===
\newcommand{\yourname}{Himanshu Kumar}%
\newcommand{\youremail}{himanshu22215@iiitd.ac.in}% <-- Update to your real email
\newcommand{\youraddress}{91-8368612369 \\ \href{https://www.linkedin.com/in/himanshu-kumar-95626634b/}{LinkedIn}}% <-- Update links
\newcommand{\yourrollno}{2022215}
\newcommand{\yourDOB}{09-03-2004}% Optional
\newcommand{\yourgithub}{https://github.com/Himanshu8368}% <-- Update
\newcommand{\yourphoto}{myphoto}% Only the filename without extension
\newcommand{\contain}{0}% Leave as is

\begin{document}
\cvheader

% === Professional Summary ===
\begin{contained}
\sectionC{Professional Summary}\newline
Computer Science student specializing in mobile application development and machine learning with proven expertise in delivering AI-powered solutions achieving 96\% accuracy using technologies like Kotlin, PyTorch, and Android development. Experienced in collaborative software development with strong foundation in data structures, algorithms, and system design, passionate about creating innovative solutions for real-world problems.
\end{contained}

% === Education Section ===
% === Education Section ===
\begin{contained}
\sectionC{Education}
\begin{longtable}{p{0.8\textwidth}p{0.19\textwidth}p{0.01\textwidth}}
    \cveducationitem{IIIT Delhi}{B.Tech (Computer Science and Engineering)}{2022 - 2026 (Present)}{CGPA: 6.33 (Till 6\textsuperscript{th} semester)}
    \cveducationitem{RSBV Patparganj, Delhi-91}{CBSE, Standard 12}{2019 - 2020}{Percentage: 84\%}
    \cveducationitem{RSBV School Block, Shakarpur, Delhi-91}{CBSE, Standard 10}{2017 - 2018}{Percentage: 82\%}
\end{longtable}
\vspace{-12pt}
\textbf{Relevant Coursework:} Operating Systems, Database Management Systems (DBMS), Algorithm Design \& Analysis, Mobile Computing, Advanced Programming
\end{contained}



% === Technical Skills Section ===
\begin{contained}
\sectionC{Technical Skills}
\begin{longtable}{p{0.25\textwidth}p{0.7\textwidth}p{0.05\textwidth}}
    \textbf{Programming Languages:}&
    Python, Java, C, C++, Kotlin\\
    \textbf{Mobile Development:}&
    Android Studio, Jetpack Compose, CameraX, Android Navigation\\
    \textbf{Machine Learning:}&
    PyTorch, PyTorch Mobile, ResNet-50, OpenAI SDK\\
    \textbf{Database \& Systems:}&
    DBMS, Operating Systems, Linux, Algorithm Design \& Analysis\\
    \textbf{Development Tools:}&
    VS Code, IntelliJ IDEA, Git, GitHub, JavaFX\\
    \textbf{Core Competencies:}&
    Data Structures \& Algorithms, System Design, Agile Development, Code Reviews\\
\end{longtable}
\end{contained}

% === Projects Section ===
% \begin{contained}
% \sectionC{Key Projects}
% \begin{longtable}{p{0.8\textwidth}p{0.19\textwidth}p{0.01\textwidth}}
%     \textbf{Plant Diseases Detection with Solution} \newline
%     \textit{Guide: Arani Bhattacharya} \newline
%     • Developed comprehensive mobile application for plant disease detection using ResNet-50 achieving 96\% accuracy on 38 disease classes \newline
%     • Implemented offline inference capabilities using PyTorch Mobile, enabling usage in remote agricultural areas without internet \newline
%     • Built intelligent advisory system using local LLMs (Llama-2/phi-2) providing contextual solutions and treatment recommendations \newline
%     • Created intuitive Android interface with Jetpack Compose, real-time camera integration, and seamless navigation \newline
%     • \textbf{Impact:} Potential to help thousands of farmers identify and treat plant diseases efficiently \newline
%     \textbf{Tech Stack:} Kotlin, Android Studio, PyTorch Mobile, OpenAI SDK \newline
%     \href{https://github.com/Himanshu8368/MC_project}{GitHub Repository} &
%     Individual Project \newline Jan 2024 - May 2025\\

%     \textbf{Stick Hero Game} \newline
%     \textit{Guide: Raghava Mutharaju} \newline
%     • Developed interactive Stick Hero game using JavaFX with smooth animations and responsive controls \newline
%     • Implemented game mechanics including physics simulation, collision detection, and score management \newline
%     • Applied object-oriented design patterns and MVC architecture for maintainable code structure \newline
%     • Collaborated with team member for UI/UX design and game logic optimization \newline
%     \href{https://github.com/Himanshu8368/Ap-Project}{GitHub Repository} &
%     Team of 2 \newline Nov 2023 – Dec 2023\\

%     \textbf{Zoo Management System} \newline
%     \textit{Guide: Raghava Mutharaju} \newline
%     • Designed comprehensive zoo management system in Java implementing advanced OOP concepts \newline
%     • Developed modules for animal management, visitor tracking, and staff scheduling with database integration \newline
%     • Implemented design patterns including Factory, Observer, and Strategy patterns for scalable architecture \newline
%     • Created command-line interface with input validation and error handling mechanisms \newline
%     \href{https://github.com/yourusername/zoo-management}{GitHub Repository} &
%     Individual Project \newline Oct 2023 – Nov 2023\\
% \end{longtable}
% \end{contained}

% === Projects Section ===
\begin{contained}
\sectionC{Projects}
\begin{longtable}{p{0.8\textwidth}p{0.19\textwidth}p{0.01\textwidth}}
    \textbf{Plant Diseases Detection with Solution} \newline
    Guide: Arani Bhattacharya \newline
    Developed AI-powered mobile application achieving 96\% accuracy in plant disease detection across 38 classes using ResNet-50. Implemented offline inference with PyTorch Mobile and integrated local LLMs for advisory system. Built Android UI with Jetpack Compose and CameraX.
    \newline \textbf{Tech Stack:} Kotlin, Android Studio, PyTorch Mobile, OpenAI SDK. \newline
    \href{https://github.com/Himanshu8368/MC_project}{GitHub} &
    Individual Project \newline Jan 2024 - May 2025\\
    \vspace{-34pt}
    \textbf{Stick Hero Game} \newline
    Guide: Raghava Mutharaju \newline
    Developed interactive Stick Hero game using JavaFX with smooth animations and responsive controls. Implemented game mechanics including physics simulation, collision detection, and applied object-oriented design patterns with MVC architecture.
    \newline \href{https://github.com/Himanshu8368/Ap-Project}{GitHub} &
    Team of 2 \newline Nov 2023 – Dec 2023\\

    \textbf{Zoo Management System} \newline
    Guide: Raghava Mutharaju \newline
    Designed comprehensive zoo management system in Java implementing advanced OOP concepts. Developed modules for animal management, visitor tracking, and staff scheduling with database integration using Factory, Observer, and Strategy patterns.
    \newline \href{https://github.com/yourusername/zoo-management}{GitHub} &
    Individual Project \newline Oct 2023 – Nov 2023\\
\end{longtable}
\end{contained}


% === Experience Section ===
\begin{contained}
\sectionC{Experience}
\begin{longtable}{p{0.8\textwidth}p{0.19\textwidth}p{0.01\textwidth}}
    \textbf{Mobile Application Developer} \newline
    \textit{Plant Disease Detection Project - IIIT Delhi} \newline
    Led development of AI-powered mobile application achieving 96\% accuracy in plant disease detection using ResNet-50 and PyTorch Mobile. Integrated local LLMs (Llama-2/phi-2) for intelligent advisory system and built responsive Android UI with Jetpack Compose. &
    Jan 2024 - May 2025\\

    \textbf{Software Development Team Member} \newline
    \textit{Multiple Course Projects - IIIT Delhi} \newline
    Collaborated in cross-functional teams of 2-4 members on software development projects using agile methodologies and version control. Contributed to full-stack development using Java, Python, and mobile technologies with focus on code quality and delivered 5+ projects on time. &
    Aug 2022 - Present\\
\end{longtable}
\end{contained}


% === Achievements Section ===
\begin{contained}
\sectionC{Achievements \& Activities}
\begin{itemize}
    \item Successfully delivered 5+ collaborative software projects with teams of 2-4 members using agile methodologies
    \item Achieved 96\% accuracy in machine learning model implementation for real-world agricultural application
    % \item Active contributor to open-source projects with well-documented GitHub repositories
    % \item Participated in multiple hackathons and coding competitions focusing on mobile and web development
\end{itemize}
\end{contained}

% === Interests Section ===
\begin{contained}
\sectionC{Interests and Hobbies}
\begin{itemize}
    \item \textbf{Sports:} Playing Table Tennis and Cricket - helps develop strategic thinking and team coordination
    \item \textbf{Geopolitics:} Following international affairs and analyzing global technological trends
    \item \textbf{Technology:} Exploring emerging technologies in AI/ML, mobile development, and open-source contributions
\end{itemize}
\end{contained}

\vspace{10pt}
% === Declaration ===
Declaration: The above information is correct to the best of my knowledge.\newline
Himanshu Kumar\newline
Date: 15/07/2025

\end{document}




% \documentclass[10pt]{extarticle}
% \usepackage{apMod}
% \usepackage{hyperref}
% % Licensed: https://www.gnu.org/licenses/gpl-3.0.html
% % Template author: Anindya Prithvi
% \vspace{-20pt}
% % === Personal Information ===
% \newcommand{\yourname}{Himanshu Kumar}%
% \newcommand{\youremail}{himanshu22215@iiitd.ac.in}% <-- Update to your real email
% \newcommand{\youraddress}{91-8368612369 \\ \href{https://www.linkedin.com/in/himanshu-kumar-95626634b/}{LinkedIn} \quad \href{https://github.com/Himanshu8368}{GitHub}}% <-- Update links
% \newcommand{\yourrollno}{2022215}
% \newcommand{\yourDOB}{09-03-2004}% Optional
% \newcommand{\yourgithub}{https://github.com/Himanshu8368}% <-- Update
% \newcommand{\yourphoto}{myphoto}% Only the filename without extension
% \newcommand{\contain}{0}% Leave as is

% \begin{document}
% \cvheader

% % === Professional Summary ===
% \begin{contained}
% \sectionC{Professional Summary}
% \begin{itemize}
%     \item Computer Science student with expertise in mobile application development and machine learning, specializing in AI-powered solutions with 96\% accuracy achievements
%     \item Experienced in collaborative software development with proven track record of delivering high-quality applications using modern technologies like Kotlin, PyTorch, and Android development
%     \item Passionate about solving real-world problems through innovative technology solutions with strong foundation in data structures, algorithms, and system design
% \end{itemize}
% \end{contained}

% % === Education Section ===
% \begin{contained}
% \sectionC{Education}
% \begin{longtable}{p{0.8\textwidth}p{0.19\textwidth}p{0.01\textwidth}}
%     \cveducationitem{IIIT Delhi}{B.Tech (Computer Science and Engineering)}{2022 - 2026 (Present)}{CGPA: 6.33 (Till 6\textsuperscript{th} semester) \newline \textbf{Relevant Coursework:} Operating Systems, DBMS, Algorithm Design \& Analysis, Mobile Computing, Advanced Programming}
%     \cveducationitem{RSBV Patparganj, Delhi-91}{CBSE, Standard 12}{2019 - 2020}{Percentage: 84\%}
%     \cveducationitem{RSBV School Block, Shakarpur, Delhi-91}{CBSE, Standard 10}{2017 - 2018}{Percentage: 82\%}
% \end{longtable}
% \end{contained}

% % === Experience Section ===
% \begin{contained}
% \sectionC{Experience}
% \begin{longtable}{p{0.8\textwidth}p{0.19\textwidth}p{0.01\textwidth}}
%     \textbf{Mobile Application Developer} \newline
%     \textit{Plant Disease Detection Project - IIIT Delhi} \newline
%     • Led development of AI-powered mobile application achieving 96\% accuracy in plant disease detection across 38 classes \newline
%     • Implemented offline inference using PyTorch Mobile, reducing response time by 80\% and enabling usage without internet connectivity \newline
%     • Integrated local LLMs (Llama-2/phi-2) for intelligent advisory system providing real-time solutions to farmers \newline
%     • Built responsive Android UI using Jetpack Compose, CameraX for image capture, and Navigation components \newline
%     \textbf{Tech Stack:} Kotlin, Android Studio, PyTorch Mobile, OpenAI SDK &
%     Jan 2024 - May 2025\\

%     \textbf{Software Development Team Member} \newline
%     \textit{Multiple Course Projects - IIIT Delhi} \newline
%     • Collaborated in cross-functional teams of 2-4 members on software development projects using agile methodologies \newline
%     • Contributed to full-stack development using Java, Python, and mobile technologies with focus on code quality and testing \newline
%     • Participated in code reviews, version control management, and delivered 5+ projects on time with high standards \newline
%     • Applied object-oriented programming principles and design patterns in game development and system design projects &
%     Aug 2022 - Present\\
% \end{longtable}
% \end{contained}

% % === Technical Skills Section ===
% \begin{contained}
% \sectionC{Technical Skills}
% \begin{longtable}{p{0.25\textwidth}p{0.7\textwidth}p{0.05\textwidth}}
%     \textbf{Programming Languages:}&
%     Python, Java, C, C++, Kotlin\\
%     \textbf{Mobile Development:}&
%     Android Studio, Jetpack Compose, CameraX, Android Navigation\\
%     \textbf{Machine Learning:}&
%     PyTorch, PyTorch Mobile, ResNet-50, OpenAI SDK\\
%     \textbf{Database \& Systems:}&
%     DBMS, Operating Systems, Linux, Algorithm Design \& Analysis\\
%     \textbf{Development Tools:}&
%     VS Code, IntelliJ IDEA, Git, GitHub, JavaFX\\
%     \textbf{Core Competencies:}&
%     Data Structures \& Algorithms, System Design, Agile Development, Code Reviews\\
% \end{longtable}
% \end{contained}

% % === Projects Section ===
% \begin{contained}
% \sectionC{Key Projects}
% \begin{longtable}{p{0.8\textwidth}p{0.19\textwidth}p{0.01\textwidth}}
%     \textbf{Plant Diseases Detection with Solution} \newline
%     \textit{Guide: Arani Bhattacharya} \newline
%     • Developed comprehensive mobile application for plant disease detection using ResNet-50 achieving 96\% accuracy on 38 disease classes \newline
%     • Implemented offline inference capabilities using PyTorch Mobile, enabling usage in remote agricultural areas without internet \newline
%     • Built intelligent advisory system using local LLMs (Llama-2/phi-2) providing contextual solutions and treatment recommendations \newline
%     • Created intuitive Android interface with Jetpack Compose, real-time camera integration, and seamless navigation \newline
%     • \textbf{Impact:} Potential to help thousands of farmers identify and treat plant diseases efficiently \newline
%     \textbf{Tech Stack:} Kotlin, Android Studio, PyTorch Mobile, OpenAI SDK \newline
%     \href{https://github.com/Himanshu8368/MC_project}{GitHub Repository} &
%     Individual Project \newline Jan 2024 - May 2025\\

%     \textbf{Stick Hero Game} \newline
%     \textit{Guide: Raghava Mutharaju} \newline
%     • Developed interactive Stick Hero game using JavaFX with smooth animations and responsive controls \newline
%     • Implemented game mechanics including physics simulation, collision detection, and score management \newline
%     • Applied object-oriented design patterns and MVC architecture for maintainable code structure \newline
%     • Collaborated with team member for UI/UX design and game logic optimization \newline
%     \href{https://github.com/Himanshu8368/Ap-Project}{GitHub Repository} &
%     Team of 2 \newline Nov 2023 – Dec 2023\\

%     \textbf{Zoo Management System} \newline
%     \textit{Guide: Raghava Mutharaju} \newline
%     • Designed comprehensive zoo management system in Java implementing advanced OOP concepts \newline
%     • Developed modules for animal management, visitor tracking, and staff scheduling with database integration \newline
%     • Implemented design patterns including Factory, Observer, and Strategy patterns for scalable architecture \newline
%     • Created command-line interface with input validation and error handling mechanisms \newline
%     \href{https://github.com/yourusername/zoo-management}{GitHub Repository} &
%     Individual Project \newline Oct 2023 – Nov 2023\\
% \end{longtable}
% \end{contained}

% % === Achievements Section ===
% \begin{contained}
% \sectionC{Achievements \& Activities}
% \begin{itemize}
%     \item Successfully delivered 5+ collaborative software projects with teams of 2-4 members using agile methodologies
%     \item Achieved 96\% accuracy in machine learning model implementation for real-world agricultural application
%     \item Implemented offline AI inference capabilities reducing application response time by 80\%
%     \item Active contributor to open-source projects with well-documented GitHub repositories
%     \item Participated in multiple hackathons and coding competitions focusing on mobile and web development
% \end{itemize}
% \end{contained}

% % === Interests Section ===
% \begin{contained}
% \sectionC{Interests and Hobbies}
% \begin{itemize}
%     \item \textbf{Sports:} Playing Table Tennis and Cricket - helps develop strategic thinking and team coordination
%     \item \textbf{Geopolitics:} Following international affairs and analyzing global technological trends
%     \item \textbf{Technology:} Exploring emerging technologies in AI/ML, mobile development, and open-source contributions
% \end{itemize}
% \end{contained}

% \vspace{10pt}
% % === Declaration ===
% Declaration: The above information is correct to the best of my knowledge.\newline
% Himanshu Kumar\newline
% Date: 15/07/2025

% \end{document}












